\section{Forelesning 3: \emph{Digitale forretningsmodeller}}
\label{sec:lec03}

\subsection*{Oversikt}
Forelesning 3 introduserer forretningsmodellen som det strategiske bindeleddet
mellom virksomhetens posisjon (\autoref{sec:lec01}) og dens digitale valg.
Hovedkilden er Weill \& Woerner (2018), \emph{What's Your Digital Business Model?},
som gir et todimensjonalt kart over digitale forretningsmodeller. Forelesningen
introduserer også Business Model Canvas (Osterwalder \& Pigneur 2010) som et
analytisk -- ikke strategisk -- verktøy.

\subsection{Hovedteori}

\subsubsection*{Forretningsmodellen som fundament}
En forretningsmodell beskriver \emph{hvordan} virksomheten skaper, leverer og
fanger verdi. Den er fundamentet som strategi bygger på -- ikke et substitutt
for strategi.

\begin{definition}{Digital forretningsmodell}
En forretningsmodell der digital teknologi er sentral for verdiskaping,
verdilevering og/eller verdifangst. Karakteriseres av mer åpne nettverk,
sterkere kunderelasjoner og dypere kundekunnskap enn klassiske modeller.
\end{definition}

\subsubsection*{Weill \& Woerner (2018): To dimensjoner}
Rammeverket plasserer virksomheter langs to akser:
\begin{enumerate}
  \item \term{Customer knowledge} -- hvor godt kjenner virksomheten sluttkunden?
    (lav $\rightarrow$ høy)
  \item \term{Business design} -- hvordan er verdinettverket organisert?
    (\emph{value chain} $\rightarrow$ \emph{ecosystem})
\end{enumerate}

Dette gir fire grunnmodeller:

\begin{definition}{De fire digitale forretningsmodellene}
\begin{enumerate}
  \item \term{Supplier} -- selger gjennom andre, lav direkte kundekunnskap,
    operer i kontrollert verdikjede. Eks.: tradisjonell merkevareleverandør.
  \item \term{Omnichannel business} -- eier kunderelasjonen på tvers av
    fysiske og digitale kanaler, moderat kundekunnskap, kontrollert kjede.
    Eks.: Coop, banker, store retailere.
  \item \term{Modular producer} -- plug-and-play-kapabiliteter solgt til
    andres økosystemer. Eks.: PayPal, Stripe; potensielt \emph{Lobyco}.
  \item \term{Ecosystem driver} -- hub i et bredt økosystem, høy
    kundekunnskap, koordinerer flere aktører. Eks.: Amazon, Alibaba, Apple.
\end{enumerate}
\end{definition}

\subsubsection*{Tre dimensjoner for konkurransefortrinn}
Weill \& Woerner identifiserer tre dimensjoner virksomheter må mestre:
\begin{itemize}
  \item \term{Content} -- hva tilbys digitalt (produkter, informasjon, tjenester).
  \item \term{Customer experience} -- hvordan digital interaksjon oppleves.
  \item \term{Platform} -- felles digital infrastruktur på tvers av aktiviteter.
\end{itemize}

\subsubsection*{Business Model Canvas (Osterwalder \& Pigneur 2010)}
BMC kartlegger forretningsmodellen i 9 byggeklosser:
\begin{itemize}
  \item Kundesegmenter, verdiforslag, kanaler, kunderelasjoner
  \item Inntektsstrømmer, nøkkelressurser, nøkkelaktiviteter
  \item Nøkkelpartnere, kostnadsstruktur
\end{itemize}

\keypoint{BMC er et analytisk verktøy, ikke en strategi i seg selv.}
Den beskriver en forretningsmodell i et snitt -- ikke konkurransedynamikk,
ikke endring over tid, ikke avveininger.

\subsection{Sentrale fagbegreper}
\begin{itemize}
  \item \term{Value proposition} -- løftet til kunden som forklarer hvorfor de
    velger oss fremfor alternativer.
  \item \term{Value creation} -- hvem som får verdi og hvordan.
  \item \term{Value capture} -- hvordan virksomheten fanger økonomisk verdi.
  \item \term{Customer knowledge} -- dybden i forståelsen av sluttkunden.
  \item \term{Ecosystem} vs.\ \term{value chain} -- nettverksbasert vs.\ lineær.
\end{itemize}

\subsection{Modeller og rammeverk}

\figureplaceholder{lec03-weill-woerner-2x2}{Weill \& Woerners 2x2-matrise:
Customer Knowledge x Business Design $\rightarrow$ fire digitale
forretningsmodeller.}

\figureplaceholder{lec03-business-model-canvas}{Business Model Canvas
(Osterwalder \& Pigneur 2010) -- 9 byggeklosser.}

\figureplaceholder{lec03-dvc-framework}{DVC Framework -- de fem dimensjonene
av digital verdi (Experiences, Relationships, Relevance, Digital Infrastructure,
Digital Outputs).}

\subsubsection*{DVC Framework (kursmateriale)}
Et komplementært rammeverk som diagnostiserer \emph{hvor} digital verdi kan
oppstå:
\begin{itemize}
  \item \term{Experiences} -- digitale interaksjoner som forbedrer kundeopplevelsen.
  \item \term{Relationships} -- digitale kanaler som utdyper kunderelasjoner.
  \item \term{Relevance} -- digital verdi tilpasset kundens situasjon.
  \item \term{Digital Infrastructure} -- back-end-kapabiliteter.
  \item \term{Digital Outputs} -- målbare digitale resultater.
\end{itemize}
\keypoint{Kritisk poeng: Digital outputs $\neq$ digital outcomes. DVC alene
kan gjøre digital aktivitet til å fremstå verdifull bare fordi den eksisterer.}

\subsection{Eksempler og caser}
\begin{itemize}
  \item \textbf{Coop} -- klassisk \emph{omnichannel business}. Etter
    lukkingen av Coop.dk MAD (2023) og Coop.dk-webshop (januar 2025) er
    posisjonen entydig store-first omnichannel.
  \item \textbf{Lobyco} -- potensiell \emph{modular producer} hvis lojalitets-
    og betalingsinfrastrukturen selges til andre retailere.
  \item \textbf{Amazon} -- prototypen på \emph{ecosystem driver}: høy
    kundekunnskap + bred koordinering av tredjeparter.
  \item \textbf{Tradisjonell merkevareprodusent} -- \emph{supplier}-modell:
    selger gjennom dagligvare uten direkte kunderelasjon.
\end{itemize}

\subsection{Eksamensrelevans}
\begin{itemize}
  \item Kunne navngi og forklare de fire DBM-typene med eksempler.
  \item Kunne plassere en konkret virksomhet (f.eks.\ Coop) i matrisen og
    forsvare valget.
  \item Forklare hvorfor BMC ikke er en strategi.
  \item Vise hvordan DVC + DBM brukes sammen (DVC: \emph{hvor} kan verdi
    skapes; DBM: blir verdien \emph{fanget} på en måte som passer modellen?).
  \item Knytte til \autoref{sec:lec04} (digital transformasjon) og
    \autoref{sec:lec06} (økosystemer).
\end{itemize}
